\documentclass[../main.tex]{subfiles}

\begin{document}
    \subsection{Kereta Api \textit{Programmer}}

    Holil ingin merantau ke jawa naik kereta. Saat tiba di stasiun, Holil diarahkan oleh petugas untuk menggunakan program sistem tiket dan gerbong kereta. Ketika program baru saja dijalankan, Holil diminta untuk memasukkan jumlah gerbong penumpang dari kereta tersebut.

    \inputfigure{res/m5i1.png}{0.6}{Keluaran Bagian 1 Program Kereta Api \textit{Programmer}}{fig:m5:i1}
    \inputfigure{res/m5i2.png}{0.3}{Keluaran Bagian 2 Program Kereta Api \textit{Programmer}}{fig:m5:i2}

    Gerbong A adalah gerbong mesin (tidak dapat diisi penumpang), sedangkan gerbong B, C, D dan seterusnya adalah gerbong penumpang. Angka di sebelah gerbong adalah jumlah penumpang yang berada di dalam gerbong tersebut.

    \inputfigure{res/m5i3.png}{0.5}{Keluaran Bagian 3 Program Kereta Api \textit{Programmer}}{fig:m5:i3}

    Kelas gerbong diatur secara otomatis oleh sistem, dimulai dari kelas EKO, BISNIS dan yang terakhir kelas RAJA. Kelas tersebut akan berulang dengan urutan yang sama sesuai jumlah gerbong.
    
    \inputfigure{res/m5i4.png}{0.43}{Keluaran Bagian 4 Program Kereta Api \textit{Programmer}}{fig:m5:i4}

    Holil dapat menambahkan penumpang baru dengan memasukkan nama, memasukkan kelas tiket, dan memilih gerbong yang tersedia.

    \inputfigure{res/m5i5.png}{0.55}{Keluaran Bagian 5 Program Kereta Api \textit{Programmer}}{fig:m5:i5}

    Holil dapat menghapus penumpang dengan memasukkan nama, memasukkan kelas tiket, dan memilih menghapus penumpang yang mana.

    \inputfigure{res/m5i6.png}{0.48}{Keluaran Bagian 6 Program Kereta Api \textit{Programmer}}{fig:m5:i6}

    Gambar di atas adalah tampilan ketika Holil sudah selesai menambah dan menghapus beberapa penumpang. Ditampilkan berapa banyak penumpang dengan kelas EKO, BISNIS dan RAJA dari seluruh gerbong kereta.

    \inputfigure{res/m5i7.png}{0.48}{Keluaran Bagian 7 Program Kereta Api \textit{Programmer}}{fig:m5:i7}

    Holil juga dapat membuang gerbong kereta mulai dari gerbong tertentu. Dalam kasus ini, Holil membuang gerbong kereta mulai dari gerbong C, sehingga gerbong C, D, E dan seterusnya sudah tidak dapat diakses lagi.
\end{document}
